\documentclass[11pt,a4paper]{article}
\usepackage{amsmath,amssymb,graphicx,geometry,hyperref,fontspec}
\geometry{margin=2.2cm}
\setlength{\parindent}{0pt}\setlength{\parskip}{4pt}
\title{LieRing Programmer Manual\\\large modules, scripts, accepted design choices}
\date{\today}
\begin{document}
\maketitle
\tableofcontents
\clearpage
\section{Programmer Manual}

Physics is in the Beam Physics Manual. This page lists
entry points of the Julia package.

The object of the ring is the Lie product of Dragt–Finn factors. That is
the same cut as MaryLie: concatenate in the algebra, analyse the product,
track a point through it. Linear optics is the closed form of
$\exp(:f_2:)$, stored as the matrix $R$. It is not a second theory of
the lattice.

\begin{figure}[htbp]\centering\includegraphics[width=\linewidth]{C:/GIT_REPOS/BELL/Julia/ring/docs/src/assets/architecture_object.png}\caption{The ring is one Lie factor: Hamiltonians are concatenated in the algebra; the product is read, tracked, or checked.}\end{figure}

\texttt{observables(lat)} builds \texttt{ring\_lie\_factor} once and reads Twiss,
$\eta$ and $\xi$ from that factor. \texttt{ring\_taylor\_map} is an RK4
polynomial and is not that object. \texttt{ring\_matrix} remains a fast
order-1 builder for plots along $s$.

\begin{figure}[htbp]\centering\includegraphics[width=\linewidth]{C:/GIT_REPOS/BELL/Julia/ring/docs/src/assets/architecture_orders.png}\caption{One polynomial, three homogeneous readings: degree 1 gives Twiss and $\eta_0$; degree 2 gives $D_2$ and $\eta_1$; degree 3 gives $\xi_1$.}\end{figure}

\subsection{Layout}

\begin{center}\small\begin{tabular}{p{0.475\linewidth}p{0.475\linewidth}}\hline
\textbf{path} & \textbf{role}\\\hline
\texttt{src/Constants.jl} & CODATA deuteron, reference particle, $v(\delta)$, $e(\delta)$, $z\leftrightarrow u$\\
\texttt{src/Jets.jl} & polynomial space, degree 4 in Hamiltonians, degree 3 in maps\\
\texttt{src/Elements.jl} & element types and lattice\\
\texttt{src/Hamiltonians.jl} & element generator $H$ as a polynomial\\
\texttt{src/LieOperators.jl} & Poisson bracket, Lie series, Dragt–Finn, product of factors to degree 4\\
\texttt{src/OptiMParser.jl} & \texttt{.opt} reader (mass from the particle, not the file header)\\
\texttt{src/LinearOptics.jl} & $R=\exp(L S \mathrm{Hess}\,H)$, slice Twiss along $s$\\
\texttt{src/Symplectic.jl} & $J^\mathsf{T}SJ-S$\\
\texttt{src/Tracking.jl} & flow of each factor, thin RF kick, product of the ring\\
\texttt{src/MapAnalysis.jl} & closed orbit, $D_2$, $\eta$, chromaticity, \texttt{observables}\\
\texttt{src/TaylorMap.jl} & independent DA map (RK4 in polynomials) — a check\\
\texttt{src/NormalForm.jl} & one-turn factor with thin RF; homological operator on generators; torus average of $\delta$\\
\texttt{src/SecondOrder.jl} & averaged longitudinal Hamiltonian, BLonD $(dt,dE)$ bridge\\
\texttt{src/ExactIntegrator.jl} & unexpanded $H$, finite cavity, renormalise\\
\texttt{src/Acceleration.jl} & ramp, tune fit\\
\texttt{src/Plotting.jl} & figures\\
\texttt{src/Optimizers.jl} & optional sextupole response\\
\hline\end{tabular}\end{center}

Entry: \texttt{include("src/LieRing.jl"); using .LieRing}.

The one-turn object used for the normal form is \texttt{one\_turn\_factor\_rf}, the
product of \texttt{ring\_lie\_factor} with the thin-kick factor. Optics of the
magnetic product without RF is \texttt{observables(lat)}. \texttt{ring\_taylor\_map}
composes truncated polynomials and is not that object.

\subsection{Scripts}

\begin{verbatim}
julia --project=. scripts/01_linear_optics.jl
julia --project=. scripts/02_lie_maps.jl
julia --project=. scripts/03_second_order.jl
julia --project=. scripts/04_optional_sextupole_optimization.jl
julia --project=. scripts/05_rf_acceleration.jl
julia --project=. scripts/06_coordinates.jl
julia --project=. scripts/07_blond_longitudinal.jl   # external drift formulae, not used in the ring product
julia --project=. scripts/08_lie_product.jl          # BCH product of two factors
julia --project=. scripts/09_lie_normal_form.jl      # one-turn factor and torus average
\end{verbatim}

The equilibrium-momentum study lives in \texttt{C:\textbackslash{}GIT\_REPOS\textbackslash{}SCT\_2\textbackslash{}Julia\textbackslash{}equilibrium\_energy}
(\texttt{scripts/run\_v08.jl} and the research note). It depends on this package
via \texttt{Pkg.develop}.

\subsection{Accepted design choices}

- Injection kinetic energy 10 MeV; final 270 MeV from the lattice header.
- Stationary study: thin RF, $V=0.2$ MV, $h=5$, $\varphi_s=\pi$.
- Ramp demo: finite cavity $L=1$ m, $\varphi_s=150^\circ$.
- Mass = CODATA deuteron $1875.61294257$ MeV (OptiM \texttt{Mass=2*938.272} ignored).
- Lattice sequence used once (file \texttt{Number of periods=2} is an error; bend is 360 deg).
- Closed $D(0)=0.011$ m (OptiM initial 0 is not the closed solution).
- Default sextupole strengths are zero.

\subsection{Docs build}

\begin{verbatim}
julia --project=docs docs/make.jl
julia docs/make_pdf.jl
julia docs/make_pdf_programmer.jl
\end{verbatim}

\end{document}
